\section{Step 8: Read the context without publishing results}
Capacity and sequence describe the caller's current state. Read the capacity
with a native doubleword load and reject values below 16\,384. Then require the
expected sequence to be below 256 and the period to be between one and 64. These
are unsigned bounds. A large value with nonzero high bits must not pass because
its low word happens to look reasonable.

Fill \texttt{record\_context}. Save the admitted sequence at stack offset 80 and
the period at offset 88. The later checksum loop will reuse temporary registers;
these stack slots retain the values it must not destroy. They belong to the
callee's private frame. No context output is written here. Rejecting an invalid
context must precede source or image access.

\paragraph{Observe} Prepare \texttt{s08-1} and stop before
\texttt{record\_context\_ready}. At \texttt{record\_context\_ready}, the eight-byte values at
\texttt{\$sp+80} and \texttt{\$sp+88} are one and four. All six context
fields still contain their entry bytes.

\Needspace{7\baselineskip}
\paragraph{Check} Where can this validation step store its temporary sequence?
\begin{itemize}
\item[A.] In the published sequence field.
\item[B.] In the private stack frame.
\item[C.] In the record checksum field.
\end{itemize}
